\documentclass[11pt,a4paper]{article}
\usepackage[utf8]{inputenc}
\usepackage[T1]{fontenc}
\usepackage{amsfonts}
\usepackage[french]{babel}
\frenchbsetup{StandardLists=true} % à inclure si on utilise \usepackage[francais]{babel}
\usepackage{enumitem}
\usepackage{amssymb}
\usepackage{amsmath}
\DeclareMathOperator{\e}{e}
\usepackage[left=2cm,right=2cm,top=2cm,bottom=2cm]{geometry}
\usepackage{multirow}
\usepackage[dvipsnames]{xcolor}

\usepackage[T1]{fontenc}
\usepackage{fouriernc}
\usepackage{graphicx}

\usepackage{setspace}
\singlespacing

\usepackage{multicol}

\usepackage{fancyhdr}
\pagestyle{fancy}
\renewcommand\headrulewidth{1pt}
\fancyhead[L]{Lycée Denis Diderot }
\fancyhead[R]{Classe de TMLEC}
\fancyfoot[C]{1}
\usepackage{multirow,tabularx,booktabs}
\newcommand{\cadreblanc}[1]{%
\medskip\framebox[\linewidth][c]{%
\rule{0mm}{#1mm}}\par
\medskip}

\setlength{\columnseprule}{.25mm}
\renewcommand{\arraystretch}{2}
\newcommand*\acocher{\quitvmode{\fboxsep0pt \fboxrule0.8pt \fbox{\vrule height1.8ex depth.1ex width0pt \vrule height0pt depth0pt width1.9ex }}}

\usepackage[tikz]{bclogo} 
\usepackage{pgfplots,mathtools}

\begin{document}
\center \LARGE{ \textbf{PROGRESSION TMELEC SCIENCES 2022 2023}}
\\[0,5cm]

\normalsize

\hrule\
\\[0,5cm]


\hrule\
\\[0,5cm]
\begin{center} \large{{Chapitre 1 - De l'alternatif au continu et inversement.} }
~~ 

~~


\hrule

\vspace{1.5 cm}
\hrule
\vspace{0.5cm}
\large{\textbf{ CCF - Lundi 17 octobre }}
~~
\vspace{0.5cm}
\hrule
\vspace{1.5cm}

\hrule
\vspace{0.5cm}
\large{{ Chapitre 2 - Réactions d'oxydo réduction et corrosion.}}
\vspace{0.5cm}    

\hrule
\vspace{1.5 cm}
\hrule
\vspace{0.5cm}
\large{{ Chapitre 3 - Propagation d'un signal sonore.}}
\vspace{0.5cm}    

\hrule
\vspace{1.5 cm}
\hrule
\vspace{0.5cm}
\large{{ Chapitre 4 - Fluide en mouvement.}}
\vspace{0.5cm}   

\hrule
\vspace{1.5 cm}
\hrule
\vspace{0.5cm}
\large{{ Chapitre 5 - Pression dans un fluide.}}
\vspace{0.5cm}   

\hrule
\vspace{1.5 cm}
\hrule
\vspace{0.5cm}
\large{{ Chapitre 6 - Moteurs électriques et puissance.}}
\vspace{0.5cm} 


\hrule
\vspace{1.5 cm}
\hrule
\vspace{0.5cm}
\large{\textbf{ CCF -  Lundi XXX mai}}
\vspace{0.5cm} 
\hrule




\end{center}

\end{document}